\chapter{Introduction and Background Research}

% You can cite chapters by using '\ref{chapter1}', where the label must
% match that given in the 'label' command, as on the next line.
\label{chapter1}

% Sections and sub-sections can be declared using \section and \subsection.
% There is also a \subsubsection, but consider carefully if you really need
% so many layers of section structure.
\section{Introduction}

Self-driving cars are becoming more popular with the aim of reducing the number of accidents related to driving. However, unlike typical computer programs that follow a set of instructions, self-driving cars must learn to drive, which involves learning to steer, brake, and navigate in order to function within unseen and complicated environments. We can simulate this using virtual environments, which allow us to simulate custom real-world environments.

An application of reinforcement learning involves learning to drive through trial and error without prior instructions. The simulated car would crash thousands of times before it understands how to perform the driving task, this approach is slow, and wastes a lot of computational power.

Methods such as imitation learning exist where we can treat the AI agent like a student, by showing examples of good driving first, instead of letting it learn through trial and error. Another way of tackling this issue is through curriculum learning, where we can gradually increase the difficulty of tasks.

My project aims to explore these teaching strategies inside a realistic simulation platform where we can train an AI agent to perform driving tasks using efficient methods.

Here is an example of how to cite a reference using the numeric system \cite{parikh1980adaptive}.


% Must provide evidence of a literature review. Use sections
% and subsections as they make sense for your project.
\section{Literature Review}

\subsection{Evolution of Autonomous Vehicles}
The implementation of autonomous vehicles initially started using a modular pipeline. As described by Chib and Singh (2024), this system involved splitting driving into different tasks: perception, localisation, planning, and control. However, this implementation was prone to issues such as error propagation, meaning that a single mistake early on within this system can cause a reaction of failures later on, ultimately leading to this method being very sensitive and leading to possibly unsafe driving.

To address these issues, end-to-end pipelines were introduced to replace these disconnected modules with a single learning pipeline. Instead of each module being trained separately, the model learns task specific features throughout the entire pipeline. This reduced the chance of failure because there is one continuous flow from input to output, so a failure in an early stage causing an error in a later stage is significantly reduced.

\subsection{Sample Inefficiency}
In reinforcement learning, an AI agent generates its own data by interacting with the environment. This leads to the cold start problem, where the agent knows nothing and takes random actions at the start of the learning process. For self-driving cars, this would involve the car driving erratically by steering, braking, and accelerating randomly.

This method of learning can be computationally expensive as driving simulators have a complex state space. The car sees a lot of information such as speed, lidar, camera, and steering angle. As the car sees all of this information in every learning episode, it processes it in all the early iterations where the car is constantly crashing and this could take millions of tries.

Many researchers aim to reduce the time and power required to teach an AI agent by reducing the number of steps required for an agent to learn.

paper to cite: https://www.cell.com/trends/cognitive-sciences/fulltext/S1364-66131930061-0

\subsection{Imitation Learning}
To address the cold start situation created by the random exploration of reinforcement learning, imitation learning observes a good existing demonstration, e.g. a demo of a car driving down a road. 

This has been used in the past for autonomous driving by using behavioural cloning. AI is fed a dataset of human driving and learns to map the camera and lidar inputs to the human's steering outputs.
**Cite paper of this being used

However, this approach also presents some challenges. The AI learns from a perfect demonstration and only learns to follow that routine, but in the event where the car makes a mistake or something different occurs in the environment, the AI doesn't know how to recover as it hasn't seen that in the demonstrations.
**Cite paper of an example

Therefore, while imitation learning is an efficient way to create a driving model, it doesn't have the ability to recover from mistakes, meaning it cannot operate safely on its own. This requires training methods that allow the agent to safely experience and recover from failures such as curriculum learning.

\subsection{Curriculum Learning}
While curriculum learning provides good initialisation for a model, curriculum learning provides a possible solution for implementing a sample efficient and robust model.




\subsection{Simulation Environment}

There are many available simulation environments to use, which are made for tasks such as reinforcement learning. They have different qualities such as high graphical fidelity for realism, or more basic ones, which require less computational power. For my research I decided to use MetaDrive as this provides a basic environment with low graphical fidelity, allowing training speeds to be significantly faster. MetaDrive also allows for the use of the OpenAI gym library, which is designed for AI tasks like reinforcement learning.

MetaDrive also provides states for the autonomous vehicle through the use of LiDar and an RGB camera, which can be used with another library such as OpenCV for object detection.

\begin{figure}[htbp]
    \centering
    \includegraphics[width=1\linewidth]{images/MetaDrive_Example.png}
    \caption{Example of the MetaDrive Simulator}
    \label{fig:placeholder}
\end{figure}